% slides/03_eeg_primer.tex
\begin{frame}{EEG 101}
  \begin{columns}[T]
    \begin{column}{0.50\textwidth}
      \textbf{Artefactos de medición}
\begin{itemize}
	\item \textit{Ocular}: parpadeos (EOG), movimientos sacádicos
	\item \textit{Muscular}: EMG ($>$30 Hz)
	\item \textit{Cardiaco}: ECG, BCG
	\item \textit{Entorno}: ruido de linea, deriva electródica
\end{itemize}
    \end{column}
    \begin{column}{0.5\textwidth}
      \vspace{0.3em}
      \begin{block}{Nota práctica}
        \small La elección de referencia (CAR, mastoides ligados, REST) impacta críticamente los resultados espectrales y de conectividad \cite{Yao2001}.
      \end{block}
    \end{column}
  \end{columns}
\end{frame}
